% ============================================================
\chapter{Espaces de Convergence et Filtres}
\label{ch:convergence}
% ============================================================

Qu'est-ce que la convergence, au fond~? Dans un cours d'analyse classique, on apprend qu'une suite $(x_n)$ converge vers $\ell$ si, pour tout voisinage de $\ell$, la suite finit par y rester. Mais cette d\'efinition pr\'esuppose l'existence d'une topologie. Dans les ann\'ees 1960, des math\'ematiciens comme Gustave Choquet, puis H.\,J.\,Kowalsky et C.\,H.\,Cook, se sont pos\'e une question audacieuse~: et si la convergence elle-m\^eme \'etait la notion primitive~? Plut\^ot que de d\'efinir d'abord les ouverts pour en d\'eduire la convergence, pourquoi ne pas sp\'ecifier directement quels filtres convergent vers quels points~? Cette inversion logique a ouvert la voie aux \emph{espaces de convergence}, une g\'en\'eralisation des espaces topologiques qui r\'esout \'el\'egamment le probl\`eme de la fermeture cart\'esienne --- un d\'efaut technique de la cat\'egorie des espaces topologiques qui compliquait l'\'etude des espaces fonctionnels.

\begin{intuition}
La notion de convergence dans un espace topologique est classiquement d\'ecrite par les
filtres ou les r\'eseaux (nets). Les espaces de convergence g\'en\'eralisent les espaces
topologiques en sp\'ecifiant directement quels filtres convergent vers quels points, sans
exiger que la structure de convergence provienne d'une topologie. Cette g\'en\'eralisation
r\'esout le probl\`eme de la fermeture cart\'esienne et est particuli\`erement adapt\'ee
aux espaces fonctionnels.
\end{intuition}

% ------------------------------------------------------------
\section{Filtres~: rappels et compl\'ements}
% ------------------------------------------------------------

\begin{definition}[Filtre]
Un \emph{filtre} sur un ensemble $X$ est une famille $\mathcal{F} \subseteq \mathcal{P}(X)$
telle que~:
\begin{enumerate}
  \item $X \in \mathcal{F}$ et $\emptyset \notin \mathcal{F}$;
  \item si $A, B \in \mathcal{F}$, alors $A \cap B \in \mathcal{F}$;
  \item si $A \in \mathcal{F}$ et $A \subseteq B$, alors $B \in \mathcal{F}$.
\end{enumerate}
\end{definition}

\begin{definition}[Base de filtre]
Une \emph{base de filtre} est une famille $\mathcal{B} \subseteq \mathcal{P}(X)$ non vide,
ne contenant pas $\emptyset$, et telle que pour tous $B_1, B_2 \in \mathcal{B}$, il existe
$B_3 \in \mathcal{B}$ avec $B_3 \subseteq B_1 \cap B_2$. Le filtre engendr\'e est
$\mathcal{F} = \{A \subseteq X : \exists B \in \mathcal{B},\, B \subseteq A\}$.
\end{definition}

\begin{definition}[Ultrafiltre]
Un \emph{ultrafiltre} est un filtre maximal (pour l'inclusion). De fa\c{c}on \'equivalente,
$\mathcal{U}$ est un ultrafiltre si pour tout $A \subseteq X$, $A \in \mathcal{U}$ ou
$X \setminus A \in \mathcal{U}$.
\end{definition}

\begin{theoreme}[Existence d'ultrafiltres]
Tout filtre est contenu dans un ultrafiltre (par le lemme de Zorn).
\end{theoreme}

\begin{definition}[Filtre image]
Si $f\colon X \to Y$ et $\mathcal{F}$ est un filtre sur $X$, le filtre image est
$f(\mathcal{F}) = \{B \subseteq Y : f^{-1}(B) \in \mathcal{F}\}$.
\end{definition}

% ------------------------------------------------------------
\section{Convergence dans les espaces topologiques}
% ------------------------------------------------------------

\begin{definition}[Filtre de voisinages]
Pour $x \in X$, le \emph{filtre de voisinages} $\mathcal{V}(x)$ est l'ensemble des
voisinages de $x$~: $\mathcal{V}(x) = \{V \subseteq X : \exists U \in \tau,\,
x \in U \subseteq V\}$.
\end{definition}

\begin{definition}[Convergence d'un filtre]
Un filtre $\mathcal{F}$ sur un espace topologique $X$ \emph{converge} vers $x$ (not\'e
$\mathcal{F} \to x$) si $\mathcal{V}(x) \subseteq \mathcal{F}$.
\end{definition}

\begin{theoreme}[Caract\'erisation de la topologie par les filtres]
Soit $X$ un espace topologique. Alors~:
\begin{enumerate}
  \item $U$ est ouvert si et seulement si pour tout $x \in U$ et tout filtre
    $\mathcal{F} \to x$, on a $U \in \mathcal{F}$.
  \item $f\colon X \to Y$ est continue si et seulement si pour tout filtre
    $\mathcal{F} \to x$, $f(\mathcal{F}) \to f(x)$.
  \item $X$ est compact si et seulement si tout ultrafiltre converge.
  \item $X$ est Hausdorff si et seulement si tout filtre a au plus une limite.
\end{enumerate}
\end{theoreme}

\begin{proof}
(3) $(\Rightarrow)$ Soit $\mathcal{U}$ un ultrafiltre. Pour chaque $x \in X$, si
$\mathcal{U} \not\to x$, il existe $U_x \in \mathcal{V}(x)$ avec $U_x \notin \mathcal{U}$,
donc $X \setminus U_x \in \mathcal{U}$. Si aucun ultrafiltre ne converge, les $U_x$
recouvrent $X$, d'o\`u un sous-recouvrement fini $U_{x_1}, \ldots, U_{x_n}$, et
$\bigcap_{i=1}^n (X \setminus U_{x_i}) = \emptyset \in \mathcal{U}$, contradiction.

$(\Leftarrow)$ Soit $(U_i)$ un recouvrement ouvert sans sous-recouvrement fini. Les
$X \setminus U_i$ ont la propri\'et\'e d'intersection finie non vide, donc engendrent un
filtre, contenu dans un ultrafiltre $\mathcal{U}$. Par hypoth\`ese, $\mathcal{U} \to x$
pour un $x$, et $x \in U_j$ pour un $j$, d'o\`u $U_j \in \mathcal{U}$ et
$X \setminus U_j \in \mathcal{U}$, contradiction.
\end{proof}

% ------------------------------------------------------------
\section{Espaces de convergence}
% ------------------------------------------------------------

\begin{definition}[Espace de convergence]
Un \emph{espace de convergence} est un couple $(X, \to)$ o\`u $X$ est un ensemble et
$\to$ est une relation entre filtres sur $X$ et points de $X$ (appel\'ee \emph{structure
de convergence}) satisfaisant~:
\begin{enumerate}
  \item Pour tout $x \in X$, le filtre principal $\dot{x} = \{A \subseteq X : x \in A\}$
    converge vers $x$~: $\dot{x} \to x$.
  \item Si $\mathcal{F} \to x$ et $\mathcal{F} \subseteq \mathcal{G}$, alors
    $\mathcal{G} \to x$.
\end{enumerate}
\end{definition}

\begin{definition}[Espace de convergence pr\'etopologique]
Un espace de convergence est \emph{pr\'etopologique} si de plus~:
\begin{enumerate}
  \setcounter{enumi}{2}
  \item Si $\mathcal{F} \to x$ et $\mathcal{G} \to x$, alors
    $\mathcal{F} \cap \mathcal{G} \to x$.
\end{enumerate}
\end{definition}

\begin{definition}[Espace de convergence topologique]
Un espace de convergence est \emph{topologique} si de plus le filtre de voisinages
$\mathcal{V}(x) = \bigcap\{\mathcal{F} : \mathcal{F} \to x\}$ converge vers $x$.
\end{definition}

\begin{theoreme}[Hi\'erarchie]
On a les inclusions strictes de cat\'egories~:
\[
  \mathbf{Top} \subsetneq \mathbf{PrTop} \subsetneq \mathbf{Conv}
\]
o\`u $\mathbf{Conv}$ est la cat\'egorie des espaces de convergence et applications
continues (pr\'eservant la convergence), $\mathbf{PrTop}$ celle des espaces
pr\'etopologiques.
\end{theoreme}

% ------------------------------------------------------------
\section{Fermeture cart\'esienne}
% ------------------------------------------------------------

\begin{theoreme}[$\mathbf{Conv}$ est cart\'esienne ferm\'ee]
La cat\'egorie $\mathbf{Conv}$ est cart\'esienne ferm\'ee~: pour tous espaces de convergence
$X$ et $Y$, l'espace $[X,Y]$ des applications continues muni de la convergence continue
est un espace de convergence, et on a~:
\[
  \mathrm{Hom}_{\mathbf{Conv}}(X \times Y, Z) \cong
  \mathrm{Hom}_{\mathbf{Conv}}(X, [Y, Z]).
\]
\end{theoreme}

\begin{definition}[Convergence continue]
Sur $[X,Y]$, un filtre $\Phi$ \emph{converge continument} vers $f$ si pour tout $x \in X$
et tout filtre $\mathcal{F} \to x$ dans $X$, le filtre \'evaluation $\Phi[\mathcal{F}]$
(engendr\'e par les $\{g(y) : g \in \phi, y \in F\}$ pour $\phi \in \Phi$,
$F \in \mathcal{F}$) converge vers $f(x)$ dans $Y$.
\end{definition}

\begin{remarque}
Le fait que $\mathbf{Conv}$ soit cart\'esienne ferm\'ee alors que $\mathbf{Top}$ ne l'est
pas est l'une des motivations principales pour \'etudier les espaces de convergence. En
s\'emantique d\'enotationnelle, cela permet de construire des espaces de fonctions sans
sortir de la cat\'egorie.
\end{remarque}

% ------------------------------------------------------------
\section{R\'eseaux (Nets) et convergence}
% ------------------------------------------------------------

\begin{definition}[R\'eseau]
Un \emph{r\'eseau} (ou \emph{net}) dans $X$ est une fonction $\varphi\colon I \to X$ o\`u
$(I, \leq)$ est un ensemble dirig\'e. On \'ecrit $(x_i)_{i\in I}$.
\end{definition}

\begin{definition}[Convergence d'un r\'eseau]
Un r\'eseau $(x_i)_{i\in I}$ converge vers $x$ dans un espace topologique si pour tout
voisinage $V$ de $x$, il existe $i_0 \in I$ tel que pour tout $i \geq i_0$, $x_i \in V$.
\end{definition}

\begin{theoreme}[\'Equivalence filtres-r\'eseaux]
Dans un espace topologique, un ensemble $F$ est ferm\'e si et seulement si toute limite
de r\'eseau dans $F$ est dans $F$. De m\^eme, $f$ est continue si et seulement si elle
pr\'eserve les limites de r\'eseaux.
\end{theoreme}

\begin{attention}
Dans les espaces non Hausdorff, un r\'eseau (ou un filtre) peut converger vers plusieurs
points \`a la fois. L'ensemble des limites d'un filtre $\mathcal{F}$ dans un espace $T_0$
est une partie basse pour l'ordre de sp\'ecialisation.
\end{attention}

% ------------------------------------------------------------
\section{Filtres de Scott et convergence}
% ------------------------------------------------------------

\begin{definition}[Filtre de Scott]
Soit $D$ un dcpo. Un filtre $\mathcal{F}$ sur $D$ est un \emph{filtre de Scott} s'il
contient un ensemble dirig\'e $S$ tel que $\bigsqcup S$ existe et $\mathcal{F}$ converge
vers $\bigsqcup S$ dans la topologie de Scott.
\end{definition}

\begin{proposition}
Dans un dcpo $D$ avec la topologie de Scott, un filtre $\mathcal{F}$ converge vers $x$ si
et seulement si $x$ est un minorant de l'ensemble des adh\'erences~:
$x \leq \bigsqcup\{\inf F : F \in \mathcal{F}, F \neq \emptyset\}$ (lorsque les infs et
sups existent). Plus pr\'ecis\'ement, $\mathcal{F} \to x$ ssi tout ouvert de Scott
contenant $x$ est dans $\mathcal{F}$.
\end{proposition}

% ------------------------------------------------------------
\section{Espaces de Choquet}
% ------------------------------------------------------------

\begin{definition}[Espace de Choquet]
Un espace topologique $X$ est un \emph{espace de Choquet} (ou espace \`a topologie
favorable) si le joueur~II a une strat\'egie gagnante dans le jeu de Choquet~: les joueurs
alternent en choisissant des ouverts non vides $U_1 \supseteq V_1 \supseteq U_2 \supseteq
V_2 \supseteq \ldots$, et le joueur~II gagne si $\bigcap_n V_n \neq \emptyset$.
\end{definition}

\begin{theoreme}
Tout espace m\'etrique complet est un espace de Choquet. Tout dcpo continu avec la
topologie de Scott est un espace de Choquet.
\end{theoreme}

% ------------------------------------------------------------
\section{Convergence et th\'eorie des domaines}
% ------------------------------------------------------------

\begin{proposition}[Convergence Scott et approximation]
Dans un domaine continu $D$, un filtre $\mathcal{F}$ converge vers $x$ dans la topologie
de Scott si et seulement si pour tout $b \ll x$, l'ensemble $\uparrow\! b \in \mathcal{F}$.
La convergence dans un domaine continu est donc d\'etermin\'ee par la relation $\ll$.
\end{proposition}

\begin{theoreme}[Espaces de convergence et domaines]
La cat\'egorie des domaines continus avec les fonctions Scott-continues est \'equivalente
\`a une sous-cat\'egorie pleine de $\mathbf{Conv}$ form\'ee d'espaces de convergence
satisfaisant des axiomes suppl\'ementaires (convergence dirig\'ee).
\end{theoreme}

% ------------------------------------------------------------
\section{Exercices}
% ------------------------------------------------------------

\begin{exercice}
Montrer qu'un espace topologique est $T_0$ si et seulement si dans la structure de
convergence associ\'ee, l'ensemble des limites de tout filtre est une partie basse
(pour l'ordre de sp\'ecialisation) et deux points distincts n'ont pas les m\^emes filtres
convergents.
\end{exercice}

\begin{exercice}
Construire un espace de convergence qui n'est pas pr\'etopologique.
\end{exercice}

\begin{exercice}
Montrer que la convergence continue sur $[X,Y]$ co\"incide avec la convergence
compacte-ouverte lorsque $X$ est localement compact de Hausdorff.
\end{exercice}

\begin{exercice}
Soit $D$ un dcpo. Montrer que le filtre engendr\'e par un ensemble dirig\'e $S$ converge
vers $\bigsqcup S$ dans la topologie de Scott.
\end{exercice}

\begin{exercice}
Montrer que dans un espace $T_0$, si un filtre $\mathcal{F}$ converge vers $x$ et vers $y$
avec $x \leq y$, alors $\mathcal{F}$ converge \'egalement vers tout $z$ avec $x \leq z
\leq y$.
\end{exercice}

\begin{exercice}[Cat\'egorie $\mathbf{Conv}$ et limites]
Montrer que $\mathbf{Conv}$ est compl\`ete et cocompl\`ete. D\'ecrire la topologie initiale
et finale en termes de convergence.
\end{exercice}
